% Standalone problem document, generated from the adjacent README.md.
% Compile with XeLaTeX. Mathematical content is maintained in Markdown.
\documentclass[11pt,a4paper]{article}
\usepackage[fontsize=10.5pt]{scrextend}
\usepackage[margin=22mm,headheight=15pt,headsep=9mm,footskip=11mm]{geometry}
\usepackage{fontspec}
\setmainfont{texgyrepagella}[Extension=.otf,UprightFont=*-regular,BoldFont=*-bold,ItalicFont=*-italic,BoldItalicFont=*-bolditalic]
\setsansfont{texgyreheros}[Extension=.otf,UprightFont=*-regular,BoldFont=*-bold,ItalicFont=*-italic,BoldItalicFont=*-bolditalic]
\setmonofont{texgyrecursor}[Extension=.otf,UprightFont=*-regular,BoldFont=*-bold,ItalicFont=*-italic,BoldItalicFont=*-bolditalic,Scale=MatchLowercase]
\usepackage{amsmath,amssymb}
\usepackage{unicode-math}
\setmathfont{texgyrepagella-math.otf}
\usepackage{xcolor,microtype,enumitem,needspace,fancyhdr}
\usepackage{bookmark}
\definecolor{ink}{HTML}{17344B}
\definecolor{accent}{HTML}{197D80}
\definecolor{muted}{HTML}{596773}
\hypersetup{colorlinks=true,linkcolor=accent,urlcolor=accent,pdftitle={FR-12 - Counting
real Hadamard
matrices},pdfauthor={Open Problems in Numerical Linear Algebra}}
\urlstyle{same}
\setlength{\parindent}{0pt}
\setlength{\parskip}{5pt plus 1pt minus 1pt}
\linespread{1.04}
\setlength{\emergencystretch}{3em}
\widowpenalty=10000
\clubpenalty=10000
\displaywidowpenalty=10000
\allowdisplaybreaks[1]
\setlist{nosep,leftmargin=1.5em}
\providecommand{\tightlist}{\setlength{\itemsep}{0pt}\setlength{\parskip}{0pt}}
\providecommand{\pandocbounded}[1]{#1}
\pagestyle{fancy}
\fancyhf{}
\fancyhead[L]{\sffamily\footnotesize\color{muted}OPEN PROBLEMS IN NUMERICAL LINEAR ALGEBRA}
\fancyhead[R]{\sffamily\bfseries\footnotesize\color{accent}FR-12}
\fancyfoot[L]{\parbox[b]{0.9\textwidth}{\sffamily\fontsize{8}{10}\selectfont\color{muted}Editorial ratings; a bounded search cannot certify that a problem remains unsolved.\\Literature check: 2026-09-12}}
\fancyfoot[R]{\sffamily\footnotesize\color{muted}\thepage}
\renewcommand{\headrulewidth}{0.3pt}
\renewcommand{\headrule}{\hbox to\headwidth{\color{accent}\leaders\hrule height \headrulewidth\hfill}}
\makeatletter
\renewcommand{\section}{\Needspace{5\baselineskip}\@startsection{section}{1}{0pt}{12pt plus 2pt minus 2pt}{4pt}{\sffamily\bfseries\large\color{ink}}}
\renewcommand{\subsection}{\Needspace{4\baselineskip}\@startsection{subsection}{2}{0pt}{9pt plus 1pt minus 1pt}{3pt}{\sffamily\bfseries\normalsize\color{ink}}}
\makeatother
\setcounter{secnumdepth}{0}
\begin{document}
{\sffamily\small\color{muted}Frames and matrix designs\par}
\vspace{3pt}
{\raggedright\hyphenpenalty=10000\exhyphenpenalty=10000\sffamily\bfseries\fontsize{21}{24}\selectfont\color{ink}Counting
real Hadamard matrices\par}
\vspace{7pt}
{\sffamily\small\textbf{Difficulty:} extreme\quad\textbf{Importance:} interesting
to the community\par}
{\sffamily\small\bfseries\color{accent}Status: Solved\par}
\vspace{4pt}
{\color{accent}\hrule height 0.8pt}
\vspace{6pt}
\textbf{Topic:} Enumeration of orthogonal sign matrices

\textbf{Rating rationale:} Known general upper bounds still have a
quadratic exponent, while the conjecture asks for an exponent of order
\(n\log n\). The question concerns the abundance of exact flat
orthogonal transforms, with connections to structured matrix
constructions and elimination.

\subsection{Negative resolution -
2026-09-12}\label{negative-resolution---2026-09-12}

\textbf{Author:} George Stepaniants, Department of Computing and
Mathematical Sciences, California Institute of Technology, Pasadena,
California, USA.

\textbf{The counting conjecture is false.} The
\href{https://github.com/ajt60gaibb/OpenProblemsInNLA/blob/main/frames-and-matrix-designs/FR-12/solution.md}{complete
proof}, Lemma 1 and Theorem 1, constructs an injection from two labeled
order-\(m\) Hadamard matrices and a perfect matching of \(2m\) row
labels. It gives

\[
H(2m)\ge(2m-1)!!\,H(m)^2,
\qquad
H(2^k)\ge 2^{\,2^k(k-1)(k-2)/8}\quad(k\ge2).
\]

The resulting exponent has order \(n(\log n)^2\) along powers of two,
contradicting every bound \(2^{C n\log_2 n}\) with fixed \(C\).
\href{https://github.com/ajt60gaibb/OpenProblemsInNLA/blob/main/frames-and-matrix-designs/FR-12/solution.pdf}{Proof
PDF} ·
\href{https://github.com/ajt60gaibb/OpenProblemsInNLA/blob/main/frames-and-matrix-designs/FR-12/solution.tex}{Standalone
proof TeX}.

The full argument passed a separate
\href{https://github.com/ajt60gaibb/OpenProblemsInNLA/blob/main/references/stepaniants-fr12-2026-09-12/REVIEW.md}{independent
Codex-agent mathematical and source-scope review}. The
\href{https://github.com/ajt60gaibb/OpenProblemsInNLA/blob/main/references/stepaniants-fr12-2026-09-12/README.md}{submission
record} preserves the supplied source, exact checks, document conversion
and bounded public fork/branch/PR and literature search. Substantial AI
assistance is disclosed; this is informal automated review, not external
human peer review or formal verification.

Ferber, Jain and Zhao retain attribution for the conjecture and prior
upper bound. This result does not settle existence at every admissible
order or determine a matching upper bound for the count. The original
statement and historical context below are retained; the difficulty and
importance ratings are historical.

\subsection{Statement}\label{statement}

For each positive integer \(n\), let

\[
H(n)=\#\{A\in\{-1,1\}^{n\times n}:AA^{\mathsf T}=nI_n\}.
\]

\textbf{Conjecture.} There is an absolute constant \(C>0\) such that

\[H(n)\le 2^{C n\log_2 n}\]

for every positive integer \(n\) divisible by four.

The count is of individual matrices with their row and column labels.
Matrices related by signed permutations are not identified. The
statement does not require that a Hadamard matrix exist at every such
order.

\subsection{Known bounds and numerical
significance}\label{known-bounds-and-numerical-significance}

Ferber, Jain and Zhao prove that some absolute \(c_H>0\) gives
\(H(n)\le 2^{(1-c_H)n^2/2}\) for every sufficiently large multiple of
four. Whenever \(H(n)>0\), distinct row permutations of one Hadamard
matrix give \(H(n)\ge n!\). Thus the conjectured exponent has the
smallest possible order along orders admitting such matrices.

Dividing a Hadamard matrix by \(\sqrt n\) produces an orthogonal
transformation whose entries all have the same magnitude. The
enumeration asks how many exact sign designs can underlie these
transforms. Peca-Medlin's work connects Hadamard enumeration for
butterfly constructions to structured orthogonal matrices and Gaussian
elimination. This is distinct from the Hadamard existence conjecture and
from bounds on elimination growth for a given matrix.

\subsection{References and status
check}\label{references-and-status-check}

\begin{itemize}
\tightlist
\item
  A. Ferber, V. Jain and Y. Zhao, \emph{On the number of Hadamard
  matrices via anti-concentration}, Combinatorics, Probability and
  Computing 31 (2022), 455--477.
  \href{https://doi.org/10.1017/S0963548321000377}{DOI};
  \href{https://www.cambridge.org/core/services/aop-cambridge-core/content/view/887EFBBF79B804BCDD942029283D4CD7/S0963548321000377a.pdf/on_the_number_of_hadamard_matrices_via_anticoncentration.pdf}{published
  PDF}. Conjecture 1.3 on p.456 is the displayed target; Theorem 1.2
  gives the upper bound. The
  \href{https://arxiv.org/abs/1808.07222}{arXiv record} also supplies
  the earlier manuscript, where the conjecture has different numbering.
\item
  \href{https://www.birs.ca/workshops/2024/24w5204/report24w5204.pdf}{BIRS
  workshop 24w5204 report (2024)}, Conjecture 28, restates the
  enumeration target.
\item
  J. Peca-Medlin, \emph{Complete pivoting growth of butterfly matrices
  and butterfly Hadamard matrices} (2026).
  \href{https://doi.org/10.1080/03081087.2026.2660796}{DOI}, §3 before
  Proposition 3.1, discusses the general count as conjectural and
  enumerates specific butterfly constructions.
\end{itemize}

On 2026-09-11, checked the published statement and later restatements,
and searched the exact title, the authors' names, Hadamard enumeration,
upper bounds, and 2025--2026 proof/counterexample combinations. No full
resolution was located. The butterfly counts apply to restricted
families and do not settle this all-matrix upper bound. This is a
bounded literature check.
\end{document}
